%======================================================================
% Appendix: Energy-Transport Audit Status for Generalized Theories
%======================================================================
% This appendix records selected inherited energy-transport claims and open
% audit items. Only the classical local Cauchy flux is derived here under
% explicit assumptions; generalized candidates are withheld or labelled open.
% No generalized family-wide law is presented as canonical.
%
% Branch: arena/40393d9f-le
% Prepared: 2026-10-21
%======================================================================

\section{Introduction}

The exact classical Eulerian flux re-derived in Section 3 is $F_i=q_i-\sigma_{ij}v_j+\mathcal{E}v_i$ (EQ-007); this appendix is a retained, unverified working collection. Unless explicitly stated otherwise, its legacy $J_i=q_i-\sigma_{ij}v_j$ expressions denote only the reduced first-order wave-flux convention, not an exact fixed-volume Eulerian total-energy flux:

\begin{equation}
J_i = q_i - \sigma_{ij}v_j,
\label{ef:classical}
\end{equation}

derived from the mechanical power balance and internal energy balance (Section 3.8 of the main manuscript).

The exact classical Eulerian flux includes advection of total energy; the reduced convention is limited to the first-order small-amplitude balance described in Section 10.2. Some generalized models may introduce additional stored-energy variables, transport channels, dissipation, or boundary-power terms; none of these is universal. This appendix retains selected source leads and open questions. It does not assert a general energy-flux law for any generalized family.

\section{Working assumptions to check model by model}

The following assumptions are provisional bookkeeping items, not a set known to hold simultaneously for every model discussed below:
\begin{enumerate}
\item Infinitesimal strain is used where a classical linearized mechanical balance is displayed.
\item Material homogeneity, isotropy, and linearity must be checked for each example.
\item Body force and heat supply are set to zero only where explicitly stated.
\item The harmonic convention $e^{-i\omega t}$ applies only to the displayed frequency-domain substitutions.
\item The classical baseline is $\mathcal{E}=\frac12\rho v_i v_i+\rho e$ with specific internal energy $e$; additional field/internal-variable storage is model-specific and not assigned here.
\item No common balance of the form $\partial_t\mathcal{E}+\partial_iJ_i=\mathcal{P}-\mathcal{D}$ is asserted for the generalized models.
\end{enumerate}

\section{Open audit item 1: Lord--Shulman (LS) Theory}

\subsection{Governing equations}

The LS theory modifies Fourier's law by introducing a flux relaxation time $\tau_0 > 0$:

\begin{equation}
q_i + \tau_0 \dot{q}_i = -k T_{,i},
\label{ef:ls:1}
\end{equation}

where $q_i$ is the heat flux, $k$ is the thermal conductivity, $T$ is the temperature increment, and $\tau_0$ is the relaxation time.

The following internal-energy balance is retained as a classical-form candidate only; its validity with the LS thermodynamic state and constitutive equations requires source checking:

\begin{equation}
\rho \dot{e} = \sigma_{ij}\dot{\varepsilon}_{ij} - q_{i,i}.
\label{ef:ls:2}
\end{equation}

The mechanical power balance (from Section 3.3) remains:

\begin{equation}
\rho \frac{d}{dt}\left(\frac{1}{2}v_i v_i\right)
= -\sigma_{ij}\dot{\varepsilon}_{ij} + (\sigma_{ij}v_j)_{,i} + \rho f_i v_i.
\label{ef:ls:3}
\end{equation}

\subsection{Classical flux form}

Adding \eqref{ef:ls:2} and \eqref{ef:ls:3}, defining $\mathcal{E} = \frac{1}{2}\rho v_i v_i + \rho e$:

\begin{equation}
\frac{\partial\mathcal{E}}{\partial t} + (q_i - \sigma_{ij}v_j + \mathcal{E}v_i)_{,i} = \rho f_i v_i.
\label{ef:ls:4}
\end{equation}

The expression $J_i=q_i-\sigma_{ij}v_j$ is only the reduced first-order wave-flux convention. The exact Eulerian flux includes $\mathcal{E}v_i$; compatibility of the reduced convention with a complete LS energy balance is not established by the constitutive substitution below and remains open.

\subsection{Frequency-domain expression}

For time-harmonic fields with $e^{-i\omega t}$ convention, \eqref{ef:ls:1} gives:

\begin{equation}
q_i(1 - i\omega\tau_0) = -k T_{,i},
\qquad
q_i = \frac{-k T_{,i}}{1 - i\omega\tau_0}.
\label{ef:ls:5}
\end{equation}

Therefore, the energy flux in the frequency domain is:

\begin{equation}
J_i = -\sigma_{ij}v_j - \frac{k T_{,i}}{1 - i\omega\tau_0}.
\label{ef:ls:6}
\end{equation}

\subsection{Time-domain interpretation}

In the time domain, \eqref{ef:ls:6} corresponds to:

\begin{equation}
J_i = q_i - \sigma_{ij}v_j,
\qquad \text{where } q_i \text{ solves } q_i + \tau_0\dot{q}_i = -k T_{,i}.
\label{ef:ls:7}
\end{equation}

There is no unique expression for $J_i$ solely in terms of $T$ in the time domain without $q_i$, because the relaxation equation introduces memory.

\subsection{Status}

\textbf{UNVERIFIED WORKING CLAIM.} Under the reduced first-order wave-flux convention only, the candidate form is $J_i = q_i - \sigma_{ij}v_j$ for LS theory with $q_i$ from \eqref{ef:ls:1}; this is not the exact Eulerian total-energy flux. The frequency-domain expression \eqref{ef:ls:6} is retained as unverified working material. A time-domain expression solely in terms of $T$ (eliminating $q_i$) would require solving the relaxation ODE, which introduces an integral (memory) term. This remains as an open formulation question.

\section{Open audit item 2: Green--Lindsay (GL) Theory}

\subsection{Status of inherited equations}

Green--Lindsay is described in the corrected manuscript as a temperature-rate-dependent, single-temperature theory with two relaxation parameters. Earlier equations in this appendix introduced a second temperature and are withdrawn as inaccurate legacy material; they are not reproduced as candidate GL equations. The complete constitutive system and its energy balance remain OPEN pending primary-source checking and derivation. The reduced first-order expression $J_i=q_i-\sigma_{ij}v_j$ is not, by itself, a verified Green--Lindsay energy-flux theorem; the exact Eulerian classical balance additionally contains total-energy advection.

\section{Open audit item 3: Green--Naghdi Type II (GN-II) Theory}

\subsection{Governing equations}

GN-II theory (thermoelasticity without energy dissipation) introduces the thermal displacement $\alpha$ defined by:

\begin{equation}
\dot{\alpha} = T,
\label{ef:gn2:1}
\end{equation}

where $T$ is the temperature increment. The source-free thermal-displacement form, obtained by integrating the primary temperature equation once with zero integration constant, is \cite{Green1993}:

\begin{equation}
\kappa^*\nabla^2\alpha = \rho c_e \ddot{\alpha} + \beta T_0 \dot{u}_{i,i}.
\label{ef:gn2:2}
\end{equation}

The exact component equation for the GN-II entropy flux has not been transcribed and verified here. No isotropic component formula is retained, and no GN-II energy-flux identity is inferred from the source's common-potential statement \cite{Green1993}.

\subsection{Energy flux}

For GN-II, there is no dissipation ($\mathcal{D} = 0$). The energy balance is:

\begin{equation}
\frac{\partial\mathcal{E}}{\partial t} + \partial_i J_i = \mathcal{P}.
\label{ef:gn2:4}
\end{equation}

The entropy flux $p_i$ is determined by the same free energy potential as the stress. The relationship between the energy flux $J_i$ and the entropy flux $p_i$ requires the full GN-II thermodynamic formulation.

\subsection{Status}

\textbf{WITHHELD.} The candidate isotropic entropy-flux expression has been removed because its exact component form is not source-verified here. The complete energy flux and its relation to entropy flux remain open pending primary-source verification and model-specific derivation.

\section{Open audit item 4: Green--Naghdi Type III (GN-III) Theory}

\subsection{Governing equations}

For the source-free, linear isotropic GN-III system transcribed from Green and Naghdi (1992), the temperature equation is:

\begin{equation}
\rho c_e\ddot{T} + \beta T_0\ddot{u}_{i,i} = \kappa\nabla^2\dot{T} + \kappa^*\nabla^2T.
\label{ef:gn3:1}
\end{equation}

\subsection{Energy flux}

The rate term can contribute dissipation when $\kappa>0$; the exact energy flux and dissipation identity are not derived here.

\subsection{Status}

\textbf{OPEN WORKING CLAIM.} The GN-III energy flux and dissipation identity require a derivation from the full 1992 formulation, including its free energy and boundary power. No mixed-flux identity is asserted here.

\section{Open audit item 5: Dual-Phase-Lag (DPL) Theory}

\subsection{Governing equations}

The DPL constitutive law (Tzou, 1995):

\begin{equation}
q_i(t + \tau_q) = -k T_{,i}(t + \tau_T),
\label{ef:dpl:1}
\end{equation}

where $\tau_q$ is the heat flux phase lag and $\tau_T$ is the temperature gradient phase lag.

Taylor expansion to first order:

\begin{equation}
q_i + \tau_q\dot{q}_i = -k(T_{,i} + \tau_T\dot{T}_{,i}).
\label{ef:dpl:2}
\end{equation}

\subsection{Energy flux}

The following internal-energy balance is retained as a classical-form candidate only; its validity with the LS thermodynamic state and constitutive equations requires source checking:

\begin{equation}
\rho\dot{e} = \sigma_{ij}\dot{\varepsilon}_{ij} - q_{i,i}.
\label{ef:dpl:3}
\end{equation}

The expression $J_i=q_i-\sigma_{ij}v_j$ is classical; its role in the displayed generalized balance remains open. However, expressing $q_i$ explicitly in terms of $T$ requires solving the DPL constitutive equation.

\subsection{Frequency-domain expression}

For time-harmonic fields:

\begin{equation}
q_i(1 - i\omega\tau_q) = -k T_{,i}(1 - i\omega\tau_T),
\label{ef:dpl:4}
\end{equation}

\begin{equation}
J_i = -\sigma_{ij}v_j - k T_{,i}\frac{1 - i\omega\tau_T}{1 - i\omega\tau_q}.
\label{ef:dpl:5}
\end{equation}

\subsection{Status}

\textbf{UNVERIFIED WORKING CLAIM.} The frequency-domain expression \eqref{ef:dpl:5} is retained as unverified working material. The time-domain expression solely in terms of $T$ requires solving the DPL constitutive equation, which introduces higher-order time derivatives. The thermodynamic admissibility condition $\tau_q \geq \tau_T$ must be verified for the flux formulation.

\section{Open audit item 6: Triple-Phase-Lag (TPL) Theory}

The inherited TPL constitutive equation and its frequency-domain expansion are withdrawn pending primary-source comparison. Different TPL truncations introduce different derivative or thermal-displacement terms, so the former formula is not retained as a working identity. The exact model, admissibility conditions, and energy balance remain OPEN; no TPL flux is derived here.

\section{Open audit item 7: Moore--Gibson--Thompson (MGT) Theory}

The dimensionally inconsistent legacy equation and its scratch frequency-domain cancellation have been withdrawn. Section~\ref{sec:4.9} records a candidate flux/thermal-displacement formulation and marks it UNVERIFIED pending a primary-source comparison. No MGT flux, stored-energy term, dissipation inequality, or frequency-domain reduction is derived here; those tasks remain OPEN.

\section{Open audit item 8: Memory-Dependent Derivative (MDD) Theory}

\subsection{Governing equations}

Candidate MDD operator (equation-level source check pending):

\begin{equation}
D_{\omega_M} f(t) = \frac{1}{\omega_M}\int_{t-\omega_M}^{t} K_{\mathrm{MDD}}(t-t') f'(t') \, dt',
\label{ef:mdd:1}
\end{equation}

where $\omega_M > 0$ is the finite memory window and $K_{\mathrm{MDD}}(t-t')$ is the kernel function.

Candidate heat-flux law coupled to the displayed operator (source check pending):

\begin{equation}
(1 + \tau_0 D_{\omega_M})q_i = -k T_{,i},
\label{ef:mdd:2}
\end{equation}

or equivalently:

\begin{equation}
q_i + \tau_0 D_{\omega_M} q_i = -k T_{,i}.
\label{ef:mdd:3}
\end{equation}

\subsection{Energy flux}

The expression $J_i=q_i-\sigma_{ij}v_j$ is classical; its role in the displayed generalized balance remains open. However, $q_i$ is now determined by a convolution-type memory equation.

\subsection{Frequency-domain expression}

For the finite-window operator in Eq.~\eqref{ef:mdd:1}, set $s=t-t'$ and use the harmonic convention $f(t)=\hat f e^{-i\Omega t}$. Its exact multiplier is kernel-dependent:

\begin{equation}
D_{\omega_M} f(t) \longleftrightarrow -\frac{i\Omega}{\omega_M}\left[\int_0^{\omega_M}K_{\mathrm{MDD}}(s)e^{i\Omega s}\,ds\right]\hat f.
\label{ef:mdd:4}
\end{equation}

This expression is the finite-horizon Fourier multiplier for the displayed kernel and convention; no kernel-independent simplified symbol is asserted.

\subsection{Status}

\textbf{OPEN WORKING CLAIM.} The MDD energy flux requires the convolution-form flux:

\begin{equation}
J_i = q_i - \sigma_{ij}v_j,
\qquad q_i + \tau_0 D_{\omega_M} q_i = -k T_{,i}.
\label{ef:mdd:5}
\end{equation}

The explicit time-domain expression for $J_i$ solely in terms of $T$ requires solving the convolution equation, which depends on the kernel choice. This derivation is open.

\section{Open audit item 9: Nonlocal Elasticity (Eringen)}

\subsection{Governing equations}

Eringen's nonlocal constitutive relation (integral form):

\begin{equation}
\sigma_{ij}(\bm{x}) = \int_\Omega \alpha(|\bm{x}' - \bm{x}|, \tau) \sigma_{ij}^{\text{local}}(\bm{x}') \, d\Omega_{\bm{x}'},
\label{ef:nonlocal:1}
\end{equation}

or in differential form (Helmholtz kernel):

\begin{equation}
(1 - \xi^2\nabla^2)\sigma_{ij} = \sigma_{ij}^{\text{local}},
\label{ef:nonlocal:2}
\end{equation}

where $\xi = e_0 a$ is the nonlocal length scale parameter.

\subsection{Energy flux}

For a specified nonlocal integral law, the local Cauchy-stress product rule need not close the domain power balance. No nonlocal energy-flux formula is asserted in this appendix. Any valid balance must be derived from the selected kernel, domain, constitutive energy, and boundary data, including whatever domain or boundary-power terms follow from those choices.

\subsection{Status}

\textbf{OPEN.} The candidate placeholder formula in the inherited working draft is intentionally withheld; it is not a derived universal correction. Derive the domain and boundary power for a specific nonlocal model before stating an energy flux or interface ratio.

\section{Open audit item 10: Piezoelectric Electrical Energy Transport}

\subsection{Governing equations}

For piezoelectric media, the constitutive relations are:

\begin{equation}
\sigma_{ij} = c_{ijkl}\varepsilon_{kl} - e_{kij}E_k,
\label{ef:piezo:1}
\end{equation}

\begin{equation}
D_i = e_{ikl}\varepsilon_{kl} + \kappa_{ij}E_j,
\label{ef:piezo:2}
\end{equation}

\begin{equation}
E_i = -\phi_{,i}, \qquad D_{i,i} = 0.
\label{ef:piezo:3}
\end{equation}

\subsection{Energy flux}

No electrical energy density or electrical power-flux expression is asserted here. The coupled energy storage and electrical work depend on the chosen electromechanical potential, field convention, external sources, and boundary conditions; the electrostatic equations displayed above alone do not determine a universal interface flux.

\subsection{Status}

\textbf{OPEN.} The inherited candidate \texttt{-phi dot(D\_i)} term and combined piezothermoelastic flux are withheld, not promoted as formulae. Derive the complete balance from a stated electromechanical free energy and field equations, then compare it with the relevant primary source.

\section{Summary Table}

\begin{table}[htbp]
\centering
\small
\begin{tabularx}{\textwidth}{@{}p{2.6cm}XX@{}}
\toprule
\textbf{Theory} & \textbf{v1.1 audit status} & \textbf{Required follow-up} \\
\midrule
Classical & DERIVED in Section 3 for local Cauchy stress and stated convention & Confirm notation consistently across all source files \\
Lord--Shulman & UNVERIFIED beyond the displayed constitutive law & Check primary-source convention, initial terms, and harmonic flux \\
Green--Lindsay & OPEN & Check the single-temperature rate-dependent equations and derive the balance \\
Green--Naghdi II/III & OPEN & Derive free energy, entropy/energy flux, and boundary power for each model \\
Dual/triple phase-lag & OPEN & Check the exact lag truncation, stability range, and energy identity \\
Moore--Gibson--Thompson & OPEN & Derive storage, boundary flux, and admissibility for the selected MGT variant \\
Memory-dependent & OPEN & Specify kernel, initial data, stored energy, and boundary power \\
Nonlocal / gradient & OPEN & Derive domain- and boundary-condition-specific power terms \\
Piezoelectric / multiphysics & OPEN & Include external-field power and derive the complete coupled balance \\
\bottomrule
\end{tabularx}
\end{table}

The preceding formulae are retained solely as unverified working material and are not approved generalized-theory identities.

\section{Recommendations}

\begin{enumerate}
\item Use only the classical flux derived in Section 3 as the canonical formula; no generalized expression in this appendix is approved.
\item Present frequency-domain forms for LS and DPL as theory-specific expressions with clearly stated assumptions.
\item Present all other formulations as open derivations with explicit assumptions.
\item Do not include any uncertain flux form in comparative tables as a canonical result.
\item Mark each uncertain flux as an open research question (OQ10.1--OQ10.8).
\end{enumerate}
